% Pursuing Stacks (cotes 134-2 to 134-8): archive cote 134-2, site page 46 (PDF page 47).
% appendix (letters to Breen). Draft compared: ps1.tex:1219-1232.
% AI-transcribed by Claude Opus 5.5 (claude-opus-5-5, adaptive thinking, effort high), Message Batches API, batch ps1 (msgbatch_01Rv5hCaNTtUt4JHx1yqoHPL), 2026-11-12; not yet checked by eye.
\scanpage{134-2/46}
\mspage{31}
classifiants au dessus de X est, je crois, développé en long et en large
dans Giraud$_{\unsure{x}}$ (il parle, si je me rappelle bien, d' "extensions" du topos X).
L'extension (si j'ose dire) de ce tapis aux n-champs ne devrait pas poser d
% NOTE: the right edge of the scan cuts the ends of some lines (here after "d")
problème.

Remords: tâchant de préciser heuristiquement la notion de topos classi
fiant d'un n-champ en groupoide\hand{s}(ou plus particulièrement, d'un n-groupoide\hand{)}
pour n\hand{$\geqslant$}2, je vois que je n'y arrive pas à vue de nez.\hand{(}Bien sûr, il suffi-%
rait (procédant de proche en proche) de savoir définir un topos classifiant
raisonnable pour une n-\uline{gerbe} , liée par un faisceau abélien \hand{$\pi$}$_n$ .\hand{)} Donc je
ne sais comment décrire le dévissage de Postnikoff en termes de n-champs,
sauf pour n \hand{$\leqslant$}2. Ceci est lié à la question d'une description directe des
groupes de cohomologie d'un n-groupoide\inserted{C}(ou d'un n-champ), à coéfficients
disons dans un système local commutatif, de façon que pour C=C$_n$(K.), K. un
ensemble semi-simplicial dont les \hand{$\pi$}$_i$ pour i\hand{$\geqslant$}n+1 sont nuls, on trouve les
groupes de cohomologie correspondants de K. . Peut-on le faire en associ-%
ant à C, de façon convenable, un ensemble semi-simplicial "nerf" de C ?
Bien entendu, si on réussit à définir un topos classifiant pour C, celui-ci
devrait être homotope à K. ci-dessus, donc avoir les mêmes invariants homo
topiques \hand{$\pi$}$_i$ et cohomologiques H$^i$; itou pour les champs. La définition habi-%
tuelle du topos classifiant, dans le cas n=1, a bien cette vertu. Cas parti
\hand{-}culier typique du Pb de la définition du topos classifiant: pour \hand{$\pi$}un grou-%
pe commutatif, trouver un topos canonique (fonctoriel en \hand{$\pi$} bien sûr ...)
ayant le type d'homotopie de K(\hand{$\pi$},n), et qui généralise la définition du
topos classifiant pour n=1 (topos des ensembles où \hand{$\pi$} opère). On frémit à
l\over{\unsure{e}}{'}idée que les topos pourraient ne pas faire l'affaire, et qu'i\unsure{l} y faille
des\hand{``}n-topos\hand{''}!! (J'espère bien que ces animaux n'existent pas ...)\margin{\unsure{(\ \ )}}.
\margin{espacer} % an arrow from below "des n-topos" goes up to the line "topos classifiant pour n=1"; another goes down from "(App.5)" to the line "ons suivantes"
\inserted{(App.5)}La théorie d' "\unsure{algèbtre|algèbre} homologique non commutative" que j'essaie de
suggérer pourrait se définir, vaguement, comme l'étude parallèle des noti-%
ons suivantes et de leurs relations multiples: a) espaces topologiques,
topos, b) ensembles semi-simpliciaux, faisceaux semi-simpliciaux etc.
c) n-catégories (notamment n-\unsure{groupoides}), n-champs (notamment n-champs en
groupoides) etc. d) complexes de groupes abéliens, de faisceaux abéliens
etc. (Les "etc" réfèrent surtout aux structures supplémentaires qu\unsure{'}on \deleted{\illegible{1}}
peut envisager sur les objets du type envisagé...) C'est donc de l'algèbr\hand{e}
avec la présence constante de motivations provenant de l'intuition topolo-%
gique. Si une telle théorie devait voir le jour, il lui faudrait bien un
nom, je me demande si "algèbre topologique" ne serait pas le plus adéquat
\hand{(}"algèbre homologique non commutative" ne peut guère aller à la longue, pour
des raisons évidentes). Ce qui est aujourd'hui parfois désigné sous ce
\hand{vocable}n'est guère qu'un bric à brac de notions \hand{(}telles que anneau topologique
\margin{46}
